[French] \documentclass[conference,twocolumn]{IEEEtran}
\IEEEoverridecommandlockouts
\usepackage{cite}
\usepackage{amsmath,amssymb,amsfonts}
\usepackage{algorithmic}
\usepackage{graphicx}
\usepackage{textcomp}
\usepackage{xcolor}
\def\BibTeX{{\rm B\kern-.05em{\sc i\kern-.025em b}\kern-.08em
    T\kern-.1667em\lower.7ex\hbox{E}\kern-.125emX}}

\begin{document}

% Title
\title{Your Paper Title Here\thanks{*This work was supported by...}}

\author{\IEEEauthorblockN{1\textsuperscript{st} Author Name}
\IEEEauthorblockA{\textit{Dept. of Your Department} \\n\textit{Your University}\\
City, Country\\
email@example.com}
\and
\IEEEauthorblockN{2\textsuperscript{nd} Author Name}
\IEEEauthorblockA{\textit{Dept. of Your Department} \\n\textit{Your University}\\
City, Country\\
email@example.com}}

\maketitle

\begin{abstract}
The abstract should briefly summarize the contents of the paper in 150--250 words.
\end{abstract}

\begin{IEEEkeywords}
component, formatting, style, styling, insert
\end{IEEEkeywords}

\section{Introduction}
This document is a model and instructions for \LaTeX\. Please observe the
conversion from program code to \LaTeX\ document. The introduction section
should provide the background and motivation for the work.

\section{Related Work}
This section should review the related work and position your contribution
in the context of existing research.

\section{Methodology}
This section should describe your methodology in detail.

\subsection{Method Overview}
Describe the overall approach and framework.

\subsection{Implementation Details}
Provide specific implementation details.

\section{Experimental Results}
This section should present and analyze the experimental results.

\subsection{Experimental Setup}
Describe the experimental environment and parameter settings.

\subsection{Results and Analysis}
Present and analyze the experimental results.

\section{Conclusion}
The conclusion should summarize the main contributions, discuss limitations,
and suggest future research directions.

\begin{thebibliography}{00}
\bibitem{b1} G. Eason, B. Noble, and I. N. Sneddon, "On certain integrals of
Lipschitz-Hankel type involving products of Bessel functions," Phil. Trans. Roy. Soc. London, vol. A247, pp. 529--551, April 1955.
\bibitem{b2} J. Clerk Maxwell, A Treatise on Electricity and Magnetism, 3rd ed., vol. 2. Oxford: Clarendon, 1892, pp.68--73.
\end{thebibliography}

\end{document}
